\section{Discussion}

This section reflects on what the implicit EM framework unifies, what it implies about standard training, and what questions remain open.

\subsection{Unification}

The framework shows that Gaussian mixture models, attention mechanisms, and cross-entropy classification are one mechanism operating under different constraints rather than three methods with superficial similarities. In GMMs, responsibilities are fully latent. In attention, they are conditioned on queries and recomputed per input. In cross-entropy, they are partially clamped by supervision. The underlying dynamics---exponentiation, normalization, responsibility-weighted updates---are identical.

This suggests a shift in how neural network training is interpreted. Probability is often treated as primitive: distributions are defined, likelihoods derived, and objectives optimized. The implicit EM perspective inverts this order. Networks compute geometric quantities, deviations from learned structures, and probabilities arise only after exponentiation and normalization. Geometry precedes probability; inference is a consequence of optimization on geometric objectives.

Loss functions, in this view, are geometric priors rather than arbitrary choices tuned for performance. Cross-entropy encodes the assumption that inputs belong to discrete categories, one of which takes full responsibility. Log-sum-exp over distances assumes inputs arise from a mixture of latent causes. Correntropy assumes outliers should be ignored. Each objective induces a different geometry of assignment and a different pattern of gradient flow. Choosing a loss function is choosing a theory of how data relates to structure.

\subsection{Implications}

For interpretability, the framework offers a direct path from training dynamics to semantic structure. If responsibilities are gradients, then the assignments a network makes are present in the backward pass, computed at every training step, rather than hidden quantities requiring probes or post-hoc analysis. The question of which component is responsible for an input has an answer in the gradient itself. This does not solve interpretability---understanding why a component takes responsibility requires further analysis---but it locates the assignment structure in a quantity that is already computed.

For objective design, the analysis reframes log-sum-exp as a structural requirement rather than a numerical convenience. Softmax is often introduced to avoid overflow or to produce well-behaved gradients, but by the analysis of Section~\ref{sec:conditions}, LSE structure is also what induces the competition that produces responsibility-weighted learning. If inference-like behavior is desired, LSE structure is required; if independent predictions or robust outlier handling are preferred, it should be deliberately avoided. The choice determines what kind of learning dynamics the objective will produce.

For theory, the framework connects two processes traditionally kept distinct: optimization concerns finding parameters that minimize a loss, while inference concerns computing posteriors over latent variables. The implicit EM result shows that under distance-based LSE objectives, these are the same process viewed at different levels: the forward pass computes posteriors and the backward pass applies them. This is not a claim that all optimization is inference; for a well-defined class of objectives, however, the distinction disappears.

\subsection{Open Directions}

Several directions remain open. The absence of volume control in most neural objectives, the missing log-determinant, leads to collapse risks currently managed by heuristics. Section~\ref{sec:volume} supplies the exact term for the linear-layer case, where the volume correction is a log-determinant of the same weights that compute the distance. What remains open is the general case: deriving implicit volume terms for deep, nonlinear compositions of distance layers, and understanding when normalization layers substitute for explicit volume control and when they do not. Answering this would connect the framework to practical stability concerns across architectures.

Supervision in real settings is rarely clean. Labels may be noisy, partial, or uncertain. The constrained regime analysis assumes hard labels that clamp responsibilities exactly; a fuller treatment would model soft or probabilistic supervision as partial constraints on the responsibility structure. This could bring semi-supervised learning, label smoothing, and learning from crowds under the same analysis.

Open-set inference requires relaxing the closed-world assumption. Current objectives force every input to be assigned; realistic deployment requires the option to reject. Objectives that support non-assignment, whether through an explicit ``none of the above'' component or a threshold below which no component takes responsibility, would extend implicit EM to settings where not all inputs belong to known categories.

Finally, the framework is falsifiable. At the loss level, the identity predicts that responsibilities can be read directly from gradients; over training, it predicts responsibility-driven specialization and, absent volume control, rich-get-richer dynamics. Tools that extract responsibilities from gradients, track specialization over training, and detect degeneration would move the framework from explanatory theory to practical instrument.
